\documentclass[10pt,a4paper]{amsart}
\usepackage{amsmath,amssymb,amsthm,graphicx}
% Retained preamble for the Stacks Project chapter divisors.tex (Divisors).
% The chapter file pinned-original-source.tex carries only \input{preamble} on its first line, so no
% preamble text lives inside the chapter itself.  The Stacks Project uses one single shared preamble for
% every chapter; the snapshot stages it verbatim in sources/stacks-chow/preamble.tex.  The block below is
% copied mechanically, line for line, from that file's lines 14-18 (the Xy-pic setup) and lines 154-217
% (the native theorem declarations together with the mathematical macros).  Nothing was added, removed or
% redefined.  The wrapper class and the amsmath/amssymb/amsthm/graphicx base packages are supplied by the
% assembler, so the class/\externaldocument part of the shared preamble is not repeated here.
% For commutative diagrams we use Xy-pic
\usepackage[all]{xy}

% We use 2cell for 2-commutative diagrams.
\xyoption{2cell}
% Theorem environments.
%
\theoremstyle{plain}
\newtheorem{theorem}[subsection]{Theorem}
\newtheorem{proposition}[subsection]{Proposition}
\newtheorem{lemma}[subsection]{Lemma}

\theoremstyle{definition}
\newtheorem{definition}[subsection]{Definition}
\newtheorem{example}[subsection]{Example}
\newtheorem{exercise}[subsection]{Exercise}
\newtheorem{situation}[subsection]{Situation}

\theoremstyle{remark}
\newtheorem{remark}[subsection]{Remark}
\newtheorem{remarks}[subsection]{Remarks}

\numberwithin{equation}{subsection}

% Macros
%
\def\lim{\mathop{\mathrm{lim}}\nolimits}
\def\colim{\mathop{\mathrm{colim}}\nolimits}
\def\Spec{\mathop{\mathrm{Spec}}}
\def\Hom{\mathop{\mathrm{Hom}}\nolimits}
\def\Ext{\mathop{\mathrm{Ext}}\nolimits}
\def\SheafHom{\mathop{\mathcal{H}\!\mathit{om}}\nolimits}
\def\SheafExt{\mathop{\mathcal{E}\!\mathit{xt}}\nolimits}
\def\Sch{\mathit{Sch}}
\def\Mor{\mathop{\mathrm{Mor}}\nolimits}
\def\Ob{\mathop{\mathrm{Ob}}\nolimits}
\def\Sh{\mathop{\mathit{Sh}}\nolimits}
\def\NL{\mathop{N\!L}\nolimits}
\def\CH{\mathop{\mathrm{CH}}\nolimits}
\def\proetale{{pro\text{-}\acute{e}tale}}
\def\etale{{\acute{e}tale}}
\def\QCoh{\mathit{QCoh}}
\def\Ker{\mathop{\mathrm{Ker}}}
\def\Im{\mathop{\mathrm{Im}}}
\def\Coker{\mathop{\mathrm{Coker}}}
\def\Coim{\mathop{\mathrm{Coim}}}

% Boxtimes
%
\DeclareMathSymbol{\boxtimes}{\mathbin}{AMSa}{"02}

%
% Macros for moduli stacks/spaces
%
\def\QCohstack{\mathcal{QC}\!\mathit{oh}}
\def\Cohstack{\mathcal{C}\!\mathit{oh}}
\def\Spacesstack{\mathcal{S}\!\mathit{paces}}
\def\Quotfunctor{\mathrm{Quot}}
\def\Hilbfunctor{\mathrm{Hilb}}
\def\Curvesstack{\mathcal{C}\!\mathit{urves}}
\def\Polarizedstack{\mathcal{P}\!\mathit{olarized}}
\def\Complexesstack{\mathcal{C}\!\mathit{omplexes}}
% \Pic is the operator that assigns to X its picard group, usage \Pic(X)
% \Picardstack_{X/B} denotes the Picard stack of X over B
% \Picardfunctor_{X/B} denotes the Picard functor of X over B
\def\Pic{\mathop{\mathrm{Pic}}\nolimits}
\def\Picardstack{\mathcal{P}\!\mathit{ic}}
\def\Picardfunctor{\mathrm{Pic}}
\def\Deformationcategory{\mathcal{D}\!\mathit{ef}}

% --- appended 2026-10-12 (Chief): the project own definitions of the reference/slogan/history
% display environments, copied mechanically from the shared Stacks preamble lines 8-12
% (\usepackage{verbatim} plus the three \newenvironment lines, which the chapter preamble keeps
% verbatim from the project preamble). Nothing added, removed or redefined: these environments
% carry no mathematical content and the project itself typesets them as comments. Without these
% three lines a counted statement whose own text uses them fails with Environment slogan undefined.
\usepackage{verbatim}
\newenvironment{reference}{\comment}{\endcomment}
\newenvironment{slogan}{\comment}{\endcomment}
\newenvironment{history}{\comment}{\endcomment}
\usepackage{cleveref}
\crefname{originalsection}{original section}{original sections}
\Crefname{originalsection}{Original section}{Original sections}
\makeatletter
\newlabel{constructions-section-projective-bundle}{{01OA}{source}{}{sourcepointer-1}{}}
\newlabel{constructions-section-projective-bundle@cref}{{[originalsection][0][]01OA}{[1][0][]source}{}{}{}}
\makeatother
\crefname{originallemma}{original lemma}{original lemmas}
\Crefname{originallemma}{Original lemma}{Original lemmas}
\makeatletter
\newlabel{lemma-relative-proj-proper}{{0800}{source}{}{sourcepointer-2}{}}
\newlabel{lemma-relative-proj-proper@cref}{{[originallemma][0][]0800}{[1][0][]source}{}{}{}}
\makeatother
\crefname{originaldefinition}{original definition}{original definitions}
\Crefname{originaldefinition}{Original definition}{Original definitions}
\makeatletter
\newlabel{morphisms-definition-modification}{{0AAZ}{source}{}{sourcepointer-3}{}}
\newlabel{morphisms-definition-modification@cref}{{[originaldefinition][0][]0AAZ}{[1][0][]source}{}{}{}}
\makeatother
\crefname{originallemma}{original lemma}{original lemmas}
\Crefname{originallemma}{Original lemma}{Original lemmas}
\makeatletter
\newlabel{morphisms-lemma-closed-immersion-proper}{{01W5}{source}{}{sourcepointer-4}{}}
\newlabel{morphisms-lemma-closed-immersion-proper@cref}{{[originallemma][0][]01W5}{[1][0][]source}{}{}{}}
\makeatother
\crefname{originallemma}{original lemma}{original lemmas}
\Crefname{originallemma}{Original lemma}{Original lemmas}
\makeatletter
\newlabel{morphisms-lemma-composition-proper}{{01W3}{source}{}{sourcepointer-5}{}}
\newlabel{morphisms-lemma-composition-proper@cref}{{[originallemma][0][]01W3}{[1][0][]source}{}{}{}}
\makeatother
\crefname{originallemma}{original lemma}{original lemmas}
\Crefname{originallemma}{Original lemma}{Original lemmas}
\makeatletter
\newlabel{morphisms-lemma-scheme-theoretic-image-reduced}{{056B}{source}{}{sourcepointer-6}{}}
\newlabel{morphisms-lemma-scheme-theoretic-image-reduced@cref}{{[originallemma][0][]056B}{[1][0][]source}{}{}{}}
\makeatother
\title{A finite locally free module on an integral scheme can be filtered after a modification}
\author{The Stacks Project Authors}
\date{stacks-project:divisors.tex:local-snapshot-20261012; GNU FDL 1.2 or later}
\begin{document}
\maketitle
\setcounter{subsection}{1}
\noindent\textit{Editorial note.} All delivered mathematical text is the author's own, copied from the pinned Stacks Project chapter file divisors.tex; no statement or proof step was rewritten, reordered or removed.\par\medskip
\noindent\textit{Editorial note.} Carrier window: statement and proof both lie inside the window 8000-10000 granted to this lane by CONTROL 8.83.\par\medskip
\noindent\textit{Editorial note.} The prefilter flagged this candidate for the phrase 'the lemma is obvious' at proof line 9422; the read shows it dismisses only the rank 0 and rank 1 cases of an induction, so the card is taken under the 2026-10-14 judge-rule (a flagged hit in a long proof is read, not dropped).\par\medskip
\noindent\textit{Editorial note.} The references of the counted spans are recorded as non-load-bearing navigation pointers; the argument that consumes them stays inside the counted proof span.\par\medskip
\noindent\textit{Editorial note.} Preamble: sources/stacks-divisors/retained-preamble.tex (88 lines, verbatim present).\par\medskip
\noindent\textit{Allowed standard prerequisites (editorial).}\begin{itemize}
\item The projective bundle and its structure morphism. Constructions, Section constructions-section-projective-bundle, used at proof lines 9425-9426. Supplies P = P(E) and the canonical surjection.
\item The relative Proj is proper. Lemma lemma-relative-proj-proper, used at proof lines 9426-9427. Makes pi proper, hence so is the restriction to X'.
\item The scheme theoretic image of a reduced scheme is integral. Morphisms, Lemma morphisms-lemma-scheme-theoretic-image-reduced, used at proof lines 9440-9442. Shows X' is integral.
\item Closed immersions and compositions of proper morphisms are proper. Morphisms, Lemmas morphisms-lemma-closed-immersion-proper and morphisms-lemma-composition-proper, used at proof lines 9442-9444. Makes f : X' -> X proper.
\item Definition of a modification. Morphisms, Definition morphisms-definition-modification, used at proof lines 9445-9446. Is the notion the constructed f is shown to satisfy.
\end{itemize}\medskip
\ifdefined\hypertarget\hypertarget{sourcepointer-1}{}\fi
\noindent\textit{Source reference.} Original section 01OA: focused-seeds/AGP/sources/stacks-snapshot-20261011/tags.txt, source lines 2187--2187. Navigation pointer into the Stacks Project; the standard fact it supplies is applied inside the counted proof span and the argument resting on it stays there.\par\medskip
\ifdefined\hypertarget\hypertarget{sourcepointer-2}{}\fi
\noindent\textit{Source reference.} Original lemma 0800: focused-seeds/AGP/sources/stacks-snapshot-20261011/tags.txt, source lines 10165--10165. Navigation pointer into the Stacks Project; the standard fact it supplies is applied inside the counted proof span and the argument resting on it stays there.\par\medskip
\ifdefined\hypertarget\hypertarget{sourcepointer-3}{}\fi
\noindent\textit{Source reference.} Original definition 0AAZ: focused-seeds/AGP/sources/stacks-snapshot-20261011/tags.txt, source lines 13002--13002. Navigation pointer into the Stacks Project; the standard fact it supplies is applied inside the counted proof span and the argument resting on it stays there.\par\medskip
\ifdefined\hypertarget\hypertarget{sourcepointer-4}{}\fi
\noindent\textit{Source reference.} Original lemma 01W5: focused-seeds/AGP/sources/stacks-snapshot-20261011/tags.txt, source lines 2470--2470. Navigation pointer into the Stacks Project; the standard fact it supplies is applied inside the counted proof span and the argument resting on it stays there.\par\medskip
\ifdefined\hypertarget\hypertarget{sourcepointer-5}{}\fi
\noindent\textit{Source reference.} Original lemma 01W3: focused-seeds/AGP/sources/stacks-snapshot-20261011/tags.txt, source lines 2468--2468. Navigation pointer into the Stacks Project; the standard fact it supplies is applied inside the counted proof span and the argument resting on it stays there.\par\medskip
\ifdefined\hypertarget\hypertarget{sourcepointer-6}{}\fi
\noindent\textit{Source reference.} Original lemma 056B: focused-seeds/AGP/sources/stacks-snapshot-20261011/tags.txt, source lines 6704--6704. Navigation pointer into the Stacks Project; the standard fact it supplies is applied inside the counted proof span and the argument resting on it stays there.\par\medskip
\begin{lemma}
\label{lemma-filter-after-modification}
Let $X$ be an integral scheme. Let $\mathcal{E}$ be a finite locally free
$\mathcal{O}_X$-module. There exists a modification $f : X' \to X$
such that $f^*\mathcal{E}$ has a filtration whose successive quotients
are invertible $\mathcal{O}_{X'}$-modules.
\end{lemma}

\begin{proof}
We prove this by induction on the rank $r$ of $\mathcal{E}$.
If $r = 1$ or $r = 0$ the lemma is obvious. Assume $r > 1$.
Let $P = \mathbf{P}(\mathcal{E})$ with structure morphism $\pi : P \to X$,
see Constructions, Section \ref{constructions-section-projective-bundle}.
Then $\pi$ is proper (Lemma \ref{lemma-relative-proj-proper}).
There is a canonical surjection
$$
\pi^*\mathcal{E} \to \mathcal{O}_P(1)
$$
whose kernel is finite locally free of rank $r - 1$.
Choose a nonempty open subscheme $U \subset X$ such that
$\mathcal{E}|_U \cong \mathcal{O}_U^{\oplus r}$.
Then $P_U = \pi^{-1}(U)$ is isomorphic to $\mathbf{P}^{r - 1}_U$.
In particular, there exists a section $s : U \to P_U$ of $\pi$.
Let $X' \subset P$ be the scheme theoretic image of the
morphism $U \to P_U \to P$. Then $X'$ is integral
(Morphisms, Lemma \ref{morphisms-lemma-scheme-theoretic-image-reduced}),
the morphism $f = \pi|_{X'} : X' \to X$ is proper (Morphisms, Lemmas
\ref{morphisms-lemma-closed-immersion-proper} and
\ref{morphisms-lemma-composition-proper}), and
$f^{-1}(U) \to U$ is an isomorphism. Hence $f$ is a modification
(Morphisms, Definition \ref{morphisms-definition-modification}).
By construction the pullback $f^*\mathcal{E}$ has a two step
filtration whose quotient is invertible because it is equal to
$\mathcal{O}_P(1)|_{X'}$ and whose sub $\mathcal{E}'$ is locally free of rank
$r - 1$. By induction we can find a modification $g : X'' \to X'$
such that $g^*\mathcal{E}'$ has a filtration as in the statement of
the lemma. Thus $f \circ g : X'' \to X$ is the required modification.
\end{proof}

\end{document}
